\section{Recipes and realization of the gradients}\label{sec:realization}

This section is deterministic.  We construct cross Gram entries
that turn a finite collection of matched vector keys into four
hole witnesses.  The construction uses the exact frame identities,
the stored nominal pin spaces and finite-dimensional linear
algebra.  It therefore applies to every frame in the support of
the thinned law $\pi$, and to every restricted leaf law, without
any conditional uniformity assumption on its channels.

We use the nominal coefficient spaces of \Cref{sec:distributions}.
For a unit $O$ and component $e$, these are
\[
 \cL_{O,e}=\bigoplus_{i\in O}(\cB_i\oplus H_i^+),\qquad
 \cR_{O,e}=\bigoplus_{i\in O}(\cB_i\oplus H_i^-),
 \qquad H_i^\pm=\F^h.
\]
Their images are obtained by summing the corresponding frame
columns in the common ambient space.  If $i'$ is the other
endpoint of $O$, put $p_i=u_{i'}U_i\in\cX^*$.

\subsection{Small tables and scalar recipes}

\begin{definition}\label{def:small-table}
For a leaf and component, put
\[
 \mathcal J_{O,e}^{\pm}
 =\bigoplus_{i\in O}(S_{i,e}^{\pm}\oplus H_i^{\pm}).
\]
For two leaves belonging to units $A,B$, a small table specifies
the two cross bilinear forms on
$\mathcal J_{A,e}^+\times\mathcal J_{B,e}^-$ and
$\mathcal J_{B,e}^+\times\mathcal J_{A,e}^-$, for every $e$.
It is admissible for two orientations if it agrees with their
actual cross Grams whenever at least one of the two table
arguments is a stored pin direction.

Let $H_{i,0}^{\pm}$ be the projection of the appropriate pin
space onto $H_i^{\pm}$, and fix a complement
$H_i^{\pm}=H_{i,0}^{\pm}\oplus H_{i,1}^{\pm}$.
The complement is called private.  We identify the two channel
coordinate copies with $\F^h$ in their common column order.
The table is injecting if, for each component, each orientation
of the roles $O,O'$, and $i\in O,z\in O'$, the following maps
are injective:
\begin{align}\label{eq:table-injection}
 D_{O,e}^+\cap\cB_i&\longrightarrow\F^h/H_{z,0}^+,
 &d&\longmapsto\text{table evaluation of $d$ against $Y_z$},\\
 D_{O,e}^-\cap\cB_i&\longrightarrow\F^h/H_{z,0}^-,
 &d&\longmapsto\text{table evaluation of $d$ against $X_z$}.
 \notag
\end{align}
These quotient conditions use the standard coordinate dot
pairing on the channels.  For instance, the first image is
zero in the quotient precisely when it pairs to zero with
all vectors annihilating $H_{z,0}^+$.
\end{definition}

Contractions on $C_i$ computed from the table carry a star.
Thus $(u_zU_i)^*(x)$, for $x\in C_i$, is computed by writing
$x_e$ in $S_{i,e}^+\otimes S_{i,e}^-$ and taking the dot product
of its two table evaluations against $Y_z,X_z$.
This is linear in $x$.  All coordinates involved have bounded
dimension, although that bound may depend on $h$.

For every cross pair $i,z$, choose between one and seven atoms
at $i$ and the same number at $z$, with a bijection preserving
the tag and the diagonal bit $q$ of each position.  At any
one endpoint all labels in its two lists must be distinct and
must avoid the exclusions of \Cref{lem:sparse-pin-labels}.
Write $v_{iz}$ for the sum of its list toward $z$.  A scalar
recipe requires, in both orientations of the two unit roles,
\begin{equation}\label{eq:recipe-scalars}
 \begin{gathered}
 \sum_{\text{list toward }z}q=1,
 \qquad a(v_{iz})+a(v_{zi})=1,\\
 T(v_{iz},\cdot)|_{C_i}
       =\bigl(a+(u_zU_i)^*\bigr)|_{C_i},\\
 (p_i+a)(v_{iz})=(p_{i'}+a)(v_{i'z}).
 \end{gathered}
\end{equation}
An atom $\mathfrak a_l(s,y)$ at $i$ has vector key
\[
 \bigl(P_{ie}v(s,y),\ Q_{ie}v(s,y)\bigr)_{e\ni l},
\]
with a fixed order of its $2(g-1)$ entries.  Matching keys means
equality in this order at every matched atom position.

\begin{lemma}[Table flags and binary prescriptions]\label{lem:small-tables}
The following statements hold.
\begin{enumerate}[label=\textup{(\roman*)}]
\item For a fixed pair of leaves and a numerical small table,
admissibility is the intersection of two unary orientation
filters, one on each unit.  Its numerical coordinate data,
including injection and the bases of the effective spaces,
have bounded description length apart from the nominal
coefficient bases, which require $O(n)$ bits.
\item The key directions of the lists at either unit are jointly
independent modulo its table-coordinate spaces, component by
component and sign by sign.
\item Given \eqref{eq:recipe-scalars}, there is a prescription of
the binary $L,R$ evaluations between distinct atoms at each
endpoint, depending only on their tags, tester summaries and
$p$ bits, which makes
\begin{equation}\label{eq:recipe-gradients}
 r(v_{iz})(x)=
 \begin{cases}
  0,&x\text{ is an atom of the list toward }z,\\
  p_{z'}(\widetilde x),&x\text{ is an atom toward }z',
 \end{cases}
\end{equation}
where $z'$ is the other endpoint of the opposite unit and
$\widetilde x$ is the matched atom there.
\end{enumerate}
\end{lemma}
\begin{proof}
A pin image is fixed on its leaf.  A cross entry with a pinned
first argument therefore depends only on the other unit's
orientation, and conversely for a pinned second argument.
Entries with both arguments pinned are numerical tests on the
leaf pair.  This proves the first assertion about filters.
Use bases of the projected pin spaces to encode tensors of
$C_i$: their coordinates lie in spaces of dimension at most
$K^2$, and there are at most $K^2$ basis tensors.  The pin,
projection and channel relationships and the table are
bounded-dimensional matrices.  The nominal bases themselves
have $O(K\dim\cB)=O(n)$ coefficients.

For the second assertion, project a purported dependence onto
an individual primal coordinate block.  It expresses a sum of
at most fourteen attack rays as an element of $S_{i,e}^{\pm}$.
Any nonzero such vector would have a sparse representation on
attack labels and is forbidden by \Cref{lem:sparse-pin-labels}.
Independence of those label factors and the nonzero constant
coordinate of each point then makes all coefficients zero.
The channels introduce no relation because the table space
is the direct sum of its individual primal and channel blocks.

For the last assertion, at one endpoint let $I_1,I_2$ be its
two disjoint nonempty lists.  The matrix of $T$ on its atoms
must be symmetric, with diagonal $a$ on each atom.
Prescribing the two list sums of its rows is possible exactly
when their self-pairings equal the corresponding diagonal
sums and their mutual pairings agree.  To see sufficiency,
subtract any symmetric matrix with the required diagonal.
The remaining task prescribes an alternating form on the
span of the two independent list indicators and its pairings
with the whole coordinate space.  Extend those indicators to
a basis, retain the prescribed entries in their rows and
columns, and assign the remaining alternating entries
arbitrarily.

For \eqref{eq:recipe-gradients}, the self-pairing requirements
hold because $T(v,v)=a(v)$.  Mutual consistency is
\[
 p_z(v_{zi})+p_{z'}(v_{z'i})
       =a(v_{iz})+a(v_{iz'}).
\]
The coupled equation at the opposite unit gives the left side
as $a(v_{zi})+a(v_{z'i})$, which equals the displayed right side
by the two role-sum equations in \eqref{eq:recipe-scalars}.
Thus the symmetric matrix exists.

Every off-diagonal entry can be implemented by formal $L,R$
bit prescriptions on that pair of atoms.  The value of $T^0$
is already determined by their tester summaries.  Choose one
indexed product term in $T^1$ linking a component of the first
star to a component of the second star; set its two factors to
give the desired correction and set all other product factors
to zero.  Forward and reverse evaluations are distinguished.
The two labels at this endpoint differ, so these are the
indexed ordered evaluations covered by
\Cref{lem:generic-mixers}.  This argument specifies the bits;
their occurrence in parameter tests will be supplied by
\Cref{lem:binary-mixing}.  It makes no assumption about
independence of arbitrary actual evaluations at fixed points.
\end{proof}

\subsection{The deterministic realization theorem}

\begin{theorem}[Gradient realization]\label{thm:gradient-realization}
Suppose two units in fixed leaves have lists satisfying
\eqref{eq:recipe-scalars}, the binary prescriptions yielding
\eqref{eq:recipe-gradients}, equal matched keys, and an
admissible injecting small table.  Freeze their actual cross
Gram entries whenever either argument is a pin or a key
direction.  Then there are abstract cross bilinear forms on
the full nominal spaces, extending all these frozen entries,
whose primal-versus-channel entries satisfy
\begin{equation}\label{eq:realized-gradients}
 u_zU_i=r(v_{iz})\quad\text{on }\cX
 \qquad(i\text{ and }z\text{ in opposite units}).
\end{equation}
Only the primal-versus-channel entries need be prescribed in
the final test.  Their number is at most
$16|\cE|h\dim\cB=O(n)$, with its coefficient fixed before
$M_0$.  Agreement of the actual cross Grams with these entries
produces all four cross holes, with witnesses $(v_{iz},v_{zi})$.
\end{theorem}
\begin{proof}
Equal keys imply
$U_iv_{iz}=U_zv_{zi}$, and the role sums in
\eqref{eq:recipe-scalars} supply the witness parity.
It remains to construct the gradients.  We distinguish
abstract Gram entries from the random event that the actual
entries agree with them.

\paragraph{A baseline of bounded primal rank.}
At each component and sign choose a splitting of the table
space as $\mathcal J=D\oplus\mathcal J'$, append the key
space, and append an additional primal complement.
The independence in \Cref{lem:small-tables} makes these
splittings possible.  Retain the small table on the product
of the two table spaces and the actual frozen entries in any
pin or key row or column.  These prescriptions agree on their
overlaps by admissibility.  Set the free blocks from an
additional primal complement to the opposite $\mathcal J'$
to zero, in both orientations, and set the additional-primal
by additional-primal block to zero.  This defines baseline
cross bilinear forms.

Fix an opposite endpoint $z$.  Let
\[
 F^0_{z,e}:\cL_{O,e}\longrightarrow\F^h,\qquad
 G^0_{z,e}:\cR_{O,e}\longrightarrow\F^h
\]
be evaluation against its $Y_z,X_z$ channels.  On additional
primal inputs, evaluation of an opposite channel factors
through its projection from $\mathcal J$ onto $D$; hence the
rank there is at most $\dim D$.  The remaining primal inputs
are in bounded-dimensional projected pin and key spaces.
It follows that the ranks of both baseline maps on all primal
inputs are bounded independently of $r_0,h,n$.  Arbitrary
channel-channel entries of the small table do not affect this
bound.  A uniform bound for either map on the combined primal
inputs of the two endpoints is
\begin{equation}\label{eq:baseline-rank-bound}
 B_{\mathrm{lin}}=3K+28:
\end{equation}
the opposite pin projection costs at most $K$ dimensions,
the two individual projected primal spaces cost at most $2K$,
and the two endpoints have at most twenty-eight key rays.

Let $u_z^0U_i$ denote the resulting contractions.  They give
the table values on $C_i$.  On an attack atom toward $z$ they
are zero, since its tensor is the equal-key tensor at $z$ and
$u_zU_z=0$.  On an attack toward $z'$ they equal
$p_{z'}$ of its matched atom, since that tensor is the
corresponding tensor at $z'$.  These are precisely the values
in \eqref{eq:recipe-gradients}.

\paragraph{Allowed changes.}
Let $\mathcal A_{O,e}^{\pm}$ be the span of all key directions
on the indicated side of the unit, and put
\[
 \mathcal F_{O,e}^{\pm}=D_{O,e}^{\pm}+\mathcal A_{O,e}^{\pm}.
\]
Write bars for the quotients
\[
 \overline{\cL}_{O,e}
 =\cL_{O,e}/\bigl(\mathcal F_{O,e}^{+}
                      +\textstyle\bigoplus_iH_{i,1}^{+}\bigr),
 \quad
 \overline{\cR}_{O,e}
 =\cR_{O,e}/\bigl(\mathcal F_{O,e}^{-}
                      +\textstyle\bigoplus_iH_{i,1}^{-}\bigr).
\]
We may change $F_z^0,G_z^0$ by any maps
\begin{equation}\label{eq:allowed-channel-changes}
 \delta F_{z,e}:\overline{\cL}_{O,e}
                     \longrightarrow(H_{z,0}^-)^\perp,
 \qquad
 \delta G_{z,e}:\overline{\cR}_{O,e}
                     \longrightarrow(H_{z,0}^+)^\perp.
\end{equation}
Here each map is extended to a bilinear change supported on
its opposite channel block and zero on the opposite primal
blocks.  It vanishes on all frozen entries: the domain
quotient kills frozen directions on one side, while on the
other side a pin's channel projection lies in $H_{z,0}$ and
is annihilated by the allowed values.  A reciprocal-role
change in the same Gram block contributes nothing on this
map's primal inputs, because it is supported on a channel
argument there.  On private-by-private blocks both changes
are zero.  They may add further channel-channel values, which
are irrelevant to the primal contractions.  Thus all choices
below may be made independently for opposite endpoints and
for the reciprocal unit roles as far as primal-input values
are concerned.

\paragraph{The unrestricted linear response.}
Fix $z$ and solve simultaneously for the two $i\in O$.
Consider the linear space of responses formed by
\[
 F_z^0\cdot\delta G_z+\delta F_z\cdot G_z^0
\]
and all bilinear forms on
$\overline{\cL}_{O,e}\times\overline{\cR}_{O,e}$, with their
contractions summed over components on each individual cut
profile.  We claim that the pair of target forms
\begin{equation}\label{eq:target-difference}
 \bigl(r(v_{iz})-u_z^0U_i\bigr)_{i\in O}
\end{equation}
belongs to this response space.

Let $(x_i)_{i\in O}\in\cX^O$ annihilate all responses.
Annihilating the arbitrary pure bilinear forms says that at
each component the sum of the two nominal tensors has zero
image in the product of barred quotients.  Project further
onto the individual primal spaces modulo
$S_{i,e}^{\pm}$ and the individual key spans.  This projection
is well defined on the barred spaces, since every pin
projects into $S_{i,e}^{\pm}$.  It shows that $x_{i,e}$ vanishes
after both such quotients.  A matrix killed after quotienting
its two factors by spaces of dimensions $d_+,d_-$ has rank
at most $d_++d_-$: in bases extending those spaces all entries
outside their rows and columns vanish.  Consequently
\[
 \rank x_{i,e}\le2K+28\le R_*.
\]
Apply \Cref{lem:label-support} to every component.

For an attack label at $i$, compare its label summand with
all key directions in that component.  Their combined label
set has at most $R_*+14$ elements, whose direction spaces
$p_s\otimes\F^{1+d_{\mathrm b}}$ form a direct sum.  Any vector of
$S_{i,e}^{\pm}$ in this sum has zero attack-label component,
by \Cref{lem:sparse-pin-labels}.  Thus the attack summand
projects to zero modulo just its individual key ray if the
atom uses this component, and is zero if there is no such
ray.  In the former case, writing its affine base vector as
$v=(1,y)$, a symmetric matrix killed in the two quotients
by $v$ has the form
\[
 Z=cvv\trans+vd\trans+dv\trans.
\]
Indeed choose $v$ as the first basis vector: the surviving
matrix can be supported only on its first row and column.
The moment identities in the original affine coordinates
now give
\[
 Z_{ii}=cv_i,\qquad Z_{0i}=cv_i+d_i+d_0v_i,
\]
so $d_i=d_0v_i$ and $Z=cvv\trans$.

The cut relation on any triangle of tags separates label by
label: the union of the three supports has at most $3R_*$
labels, within the interpolation bound.  At this attack
label the nonzero components lie on its tag star; the
triangle relation forces their scalar coefficients to be
equal.  Subtract that scalar multiple of the attack atom
from $x_i$.  Do this for every attack label.  Each subtracted
atom annihilates the response space, because both its factor
directions are keys.  The remaining profiles still
annihilate the response and have only nonattack label
support.  Removing these separated label blocks does not
increase their component ranks or label support sizes.

At one component contract a remaining $x_{i,e}$ on the minus
side by a primal functional $\phi$ at $i$ annihilating
$S_{i,e}^-$ and all individual key directions.  Extend $\phi$
by zero on other endpoint and channel blocks; it is a
functional on $\overline{\cR}_{O,e}$.  Pure annihilation shows
that the resulting individual plus vector $d$ is a sum of a
pin, key directions, and private channel directions.  In
this representation all key coefficients vanish.  To see
this, project onto each individual primal block.  At the
other endpoint the projected pin would be a sparse sum of
attack rays.  At endpoint $i$ it is the difference of the
nonattack vector $d$, supported on at most $R_*$ labels, and
a sum of at most fourteen attack rays.  In both cases the
label exclusion prevents any attack coefficient.  The
channel components of the pin lie in the protected spaces
$H_0$, whereas the remaining channel term lies in their
chosen complements $H_1$.  They must both vanish.  Therefore
\[
 d\in D_{O,e}^+\cap\cB_i.
\]
Test the derivative response using
$\delta G_{z,e}=\phi\otimes c$ for arbitrary
$c\in(H_{z,0}^+)^\perp$.  It gives
$F_{z,e}^0(d)\cdot c=0$.  Hence the table evaluation of $d$
is zero modulo $H_{z,0}^+$, and injection in
\eqref{eq:table-injection} yields $d=0$.
The reciprocal contraction proves the same assertion in the
minus mode.

It follows that the remaining component tensors lie in the
products of the individual projected pin plus key spaces.
Their nonattack support removes the key directions by the
same sparse-label argument.  Thus the remaining profiles
belong to $C_i$.  We have proved that every annihilator is a
sum of effective-space profiles and attack atoms.  The target
\eqref{eq:target-difference} vanishes on the former by
\eqref{eq:recipe-scalars} and on the latter by
\eqref{eq:recipe-gradients} and the baseline values.
In a finite-dimensional vector space the image of a linear
map is the annihilator of the kernel of its dual: this
follows by extending a basis of the image to a basis of the
target space.  It proves the claimed linear solvability.

\paragraph{A low-rank representation of the target.}
For one $i$, the form $r(v_{iz})$ has a matrix representation
on the components of rank at most
\begin{equation}\label{eq:unconstrained-gradient-rank}
 15r_0+14J|\cE|
\end{equation}
per component.  Here is an explicit routing for its tester
part.  Let $w=(w_t)$ be the representation given by its atom
list.  The list condition gives $\chi_*(w)=\sum q=1$.
On an allowed $O_{d,t}$ block at any tag the coefficient in
$a+T^0(v_{iz},\cdot)$ is
\[
 1+a(v_{iz})+b_t(v_{iz})
 =1+\chi_*(w)+\eta_S(w_t)=\eta_S(w_t).
\]
There are at most seven tags $t$ with nonzero $\eta_S(w_t)$.
Route this tester from each allowed tag $l\in I_d$ to $d$
through the component $\{l,d\}$.  It contributes nothing at
$d$, where $O_{d,t}$ is disallowed.  At a component this uses
at most fourteen ordinary testers, allowing either endpoint
as center.  For the shared tester, route a star centered at
$t$ with coefficient $a_t(v_{iz})$.  Its contribution at the
center is $(g-1)a_t(v_{iz})=0$, and at any other tag it gives
the desired term in $a_t(v_{iz})b_t$.  Consolidating gives at
most one shared tester per component.  Each ordered tester
matrix has rank $r_0$, proving the $15r_0$ bound.

For $T^1(v_{iz},\cdot)$, each of the two orientations of an
indexed term supplies a matrix of rank at most the component
rank of $v_{iz}$.  This profile is a sum of at most seven
atoms, so its total component rank is at most $7(g-1)$.
The total contribution per component has rank at most
$14J(g-1)\le14J|\cE|$.  This proves
\eqref{eq:unconstrained-gradient-rank}.

\paragraph{Removing the bounded-dimensional obstructions.}
Let $M_{i,e}$ be the preceding representing matrices.  On each
individual primal factor choose a projection $A_{i,e}^{\pm}$
whose kernel is the sum of its projected pin and key spaces.
The matrix
$(A_{i,e}^+)\trans M_{i,e}A_{i,e}^-$, extended by zero outside
that individual block, is a permitted pure quotient form.
Its sum over $i$ has rank at most twice
\eqref{eq:unconstrained-gradient-rank}.  Moreover
\begin{align}\label{eq:bounded-rank-removal}
 M_{i,e}-(A_{i,e}^+)\trans M_{i,e}A_{i,e}^-
 &= (I-A_{i,e}^+)\trans M_{i,e}
       +(A_{i,e}^+)\trans M_{i,e}(I-A_{i,e}^-),\\
 \rank\bigl(M_{i,e}-(A_{i,e}^+)\trans M_{i,e}A_{i,e}^-\bigr)
 &\le \rank(I-A_{i,e}^+)+\rank(I-A_{i,e}^-)
 \le 2K+28.\notag
\end{align}
Subtract these pure forms from the linear task.  Its residual,
including the baseline contraction, has a per-individual
matrix representation of bounded rank independent of
$r_0,h,n$, and remains linearly soluble.

In any unrestricted solution, replace the derivative output
maps by bounded-rank maps in the same allowed value spaces,
preserving all their dot products with the opposite baseline
values on primal inputs.  Such replacement is possible as
follows.  The map from the allowed value space to the dual
of the span of those baseline values has bounded rank.
Choose a linear section of its image and compose with that
map.  This has bounded-dimensional image and preserves every
observed pairing.  Postcomposing the original derivative map
with this value projection leaves its primal derivative
response unchanged and preserves its domain quotient.
The two new derivative maps have bounded rank.  Their extra
product $\delta F_z\cdot\delta G_z$ is also a bounded-rank
pure form, which we include with opposite sign in the
remaining pure task.  That task still has an unrestricted
pure solution.  Its per-individual target matrices have rank at
most
\begin{equation}\label{eq:residual-rank-bound}
 R_{\mathrm{res}}=2K+28+4B_{\mathrm{lin}}=14K+140.
\end{equation}
Here the first term bounds the projection difference, one
$B_{\mathrm{lin}}$ bounds the baseline contraction, two bound
the derivative terms, and one bounds their extra product.

\paragraph{Compression of the remaining pure task.}
We show that the latter task has a pure solution of bounded
rank.  At endpoint $i$, expand all projected pin and key
vectors, and all row and column forms in its bounded-rank
target representation, along selector coordinates and base
blocks.  Their counts depend on earlier construction
parameters but not on $r_0,h,n$.  On each ordinary base block
let $U$ be the span of the vectors to fix and $K_0$ the common
kernel of the forms to preserve.  Choose a complement in
$K_0$ to $U\cap K_0$, and project along that complement onto
a subspace containing $U$.  The resulting endomorphism fixes
$U$, preserves all the forms, and has rank at most
$\dim U+\codim K_0$.  There are at most
$2|\cE|(K+14)$ vectors to fix and at most
$2|\cE|R_{\mathrm{res}}$ forms to preserve at one endpoint
before selector expansion.  Thus an explicit bound for the
rank on an ordinary block is
\begin{equation}\label{eq:block-compression-bound}
 D_{\mathrm{blk}}
 =2p_*|\cE|(K+14+R_{\mathrm{res}}).
\end{equation}

Use these endomorphisms blockwise, fix the constant base
coordinate, and act trivially on the selector factor.  Use
the same resulting map at every component and in both modes
at this endpoint.  It sends each allowed point to an allowed
point and therefore preserves all $W_l$ and the cut domain.
It fixes the individual projected pin and key vectors.
Taking identity on channels consequently fixes every mixed
pin itself and preserves private channel subspaces.  The
map descends to the barred spaces and has bounded rank there:
only the bounded protected channel dimensions survive their
quotients.  On either barred space a uniform bound is
\begin{equation}\label{eq:quotient-compression-bound}
 D_{\mathrm{quo}}
 =2p_*\bigl(1+(g^2+3)D_{\mathrm{blk}}\bigr)+2K.
\end{equation}
The first term counts the two compressed primal blocks and
the second bounds their surviving channel projections.

Precompose an unrestricted pure solution on its two arguments
with these maps.  Its rank per component becomes bounded.
On individual cut profiles its value is unchanged, because
the compressed profile remains in that cut domain and every
form in the target representation is preserved.  It is
therefore the required bounded-rank pure solution.

\paragraph{Factoring the pure forms through channels.}
Restore the earlier projected pure forms.  Their sum with
the compressed correction has rank at most
\begin{equation}\label{eq:final-pure-rank}
 30r_0+28J|\cE|+D_{\mathrm{quo}}
\end{equation}
per component for the fixed opposite endpoint.  In particular
the additive constant is independent of $r_0,h,n$.
In the two allowed channel value spaces impose, in addition,
orthogonality to the opposite baseline and bounded derivative
values on primal inputs.  These are only boundedly many
linear restrictions.  Each remaining value space has
codimension at most $K+2B_{\mathrm{lin}}$, so their dot pairing
has rank at least $h-2K-4B_{\mathrm{lin}}$: restricting a
nondegenerate pairing by codimensions $c_1,c_2$ loses at most
$c_1+c_2$ in rank.  With $h=1000r_0$, the sufficient requirement
is
\begin{equation}\label{eq:realization-r0-requirement}
 970r_0>28J|\cE|+D_{\mathrm{quo}}+2K+4B_{\mathrm{lin}}.
\end{equation}
Every quantity on its right is fixed before $r_0$, so this is
one of the permitted requirements on that choice.

Factor a rank-$r$ desired pure bilinear form as
$\sum_{a=1}^r\ell_a\otimes m_a$.  Choose vectors $f_a,g_a$
in these remaining channel spaces with $f_a\cdot g_b$ equal
to the Kronecker delta; they exist by the pairing rank.
The changes $\sum_a\ell_af_a$ and $\sum_am_ag_a$ realize
that form as their dot product.  Their cross terms with the
baseline and the previous small derivative changes vanish
on primal inputs by the additional orthogonality conditions.
They obey \eqref{eq:allowed-channel-changes} and hence
preserve every frozen entry.  This realizes the target
exactly for both endpoints $i\in O$.

Repeat for both opposite endpoints and both unit roles.
The independence of primal-input changes already proved
ensures that the resulting two cross Gram blocks realize
all equations \eqref{eq:realized-gradients} simultaneously.
Only primal-channel entries have been used.  In each of the
two blocks there are at most $8h\dim\cB$ such entries per
component, proving the asserted count.  Actual agreement
with them gives all gradient equations; equal keys and the
role sums give the other witness conditions.  All four
cross pairs are therefore holes.
\end{proof}
